\scanpage{6}
\mspage{3}
\uline{Rigidité}: $M\longrightarrow M^\vee$
\inserted{\deleted{\unsure{$=D(M)$}} $D(M)$}
compatible avec foncteur fibre $F'$,
foncteur équivalence de catégories, donc \uline{exact}
donc cont. à droit, \deleted{\unsure{donc donné aussi}}
donc donné par
\[
M\longmapsto\operatorname{Hom}_A(M,I)
\]
$I$ objet de $C$, muni d'opération de $A$: g.
\[
I=\deleted{\unsure{U}}D(A)
\underset{\inserted{\text{compatibilité avec }F'}}{=}
A^{\vee(k')}=U
\]
avec opération à g. de $A$ provenant par transposition
des opérations de $A$ à droite sur $A_g$, donc $U$ est un $A$-module.

Il faut donc faire opérer $A$ à gauche sur cet objet de
$\operatorname{Mod}(A)$ (donc mettre sur $U$ une structure de
\deleted{bimodule} module sur
\[
\underbrace{A}_{\text{op. à définir}}\otimes
\underbrace{A}_{\text{opération déjà connue}}.
\]

\margin{\uline{NB}}

Or les endom de $U$ en tant que $A$-module à g.
\inserted{(attention, ils ne sont pas néc. $k'$-linéaires,
pour la $k'$-structure \unsure{donnée|dualité} sur $U$!)}
correspondent aux endomorphismes de $A$ en tant que $A$ module à droite:
\[
\operatorname{End}_A(U)\simeq
\operatorname{End}_A(A_d)^o\simeq A^o
\]
donc il faut se donner
\[
A\xrightarrow{a:u\mapsto\unsure{\check u|\widetilde u}}A^o
\]

\margin{\unsure{est-ce que} \unsure{n'est valable que|n'est valable} \unsure{si}
$k'/k$ \unsure{est fini}, \unsure{donc} pour identifier
$U$ à $D_k(A)$ $(\simeq D_{k'}(A))$}
\margin{5}
